\chapter{Estado del arte} \label{cap2}

\vspace*{5mm}

%Introducción%
En esta sección explicaremos los elementos necesarios para la resolución del problema propuesto. Haremos una explicación en profundidad de los sistemas de aprendizaje federados y de la clasificación multi-etiqueta.

\section{Aprendizaje Federado}
En aprendizaje automático tenemos un modelo que entrenamos conuna serie de datos para que realize una tarea que puede ser reconcoer objetos en una imagen, controlar la reacción de un sistema ante una situación o jugar a un juego. En la práctica estos datos no se generan en la máquina dónde entrenamos al modelo sino que pueden venir de un dispositivo móvil, unos sensores o un micrófono. Estos datos también pueden venir no de una única fuente sino de varios lugares que recopilen información para realizar una tarea en conjunto. La aproximación más común ante esta estructura de datos es unir todas las fuentes en un lugar centralizado y tratarlas para crear un único conjunto de datos y usarlo para realizar el entrenamiento de nuestro modelo.

Un entorno de aprendizaje federado tiene como objetivo el entrenamiento de un modelo de aprendizaje automático mediante la colaboración de distintos clientes. Los datos están también federados de manera que cada cliente contiene un subconjunto de los datos totales de manera que permanecen decentralizados. Cada cliente entrena de manera local su modelo y una vez finaliza el entrenamiento comparte los pesos del modelo entrenado con el resto de clientes de manera que se agregan para actualizar los modelos locales con todos los entrenamientos de la red, de esta forma no se comparten los datos de los clientes pero el conjunto de clientes se puede aprovechar del entrenamiento que han realizado los dema´s.

\section{Clasificación multi-etiqueta}

El estudio Heterogeneous Federated Learning: State-of-the-art and Research Challenges ~\cite{Ye:HFL_Survey} realiza una clasificación de los problemas presentes en las diferentes arquitecturas de modelos federados según si la heterogeneidad es estadística, en el modelo, en las comunicaciones o en los dispositivos. Si hablamos de la heterogeneidad estadística estamos hablando de los datos y el principal problema es que no sean independientes e identicamente distribuídos (no-IID) lo que hace que el aprendizaje se reduzca considerablemente ~\cite{Zhao_ML}~\cite{criado2021noniiddatacontinuallearning}. Ahondando más en esta categorización pueden tener desbalance de etiquetas, desbalance de características, desbalance de calidad y desbalance de cantidad. Artículos como ~\cite{Zhaobin:FedMLP} hablan de este problema en datos médicos y atacan el problema centrándose en la normalización de los datos.

El estudio también evalúa otros tipos de heterogeneidad en la comunicación, los dispositivos y los modelos y propone diferentes soluciones para mitigar estos problemas. El caso que nos compete tiene desbalance en etiquetas ya que como hemos comentado antes cada cliente va a entrenar una tarea distinta, por lo que los datos de cada uno solo tendrán un subconjunto de las etiquetas y también desbalance en los modelos ya que cada cliente va a tener un número de etiquetas diferente por lo que la salida de la red local no será consistente entre todo el sistema. En cuanto al desbalance en el modelo se proponen varias soluciones, una de ellas es la compartición de arquitectura, que consiste en compartir parte del modelo y hacer un entrenamiento personalizado, este método es capaz de solucionar tanto la heterogeneidad estadística como la de modelo y de dispositivo al mismo tiempo ~\cite{Manoj:FedPer}.

Los métodos a nivel de datos pueden resolver el problema asociado a las grandes diferencias en la distribución de los datos y acelerar la convergencia reduciendo heterogeneidad estadística entre ellos. Por el contrario, los métodos a nivel de modelo buscan aprender un modelo local para cada cliente que se adapte a su distribución de datos privada, esta solución es la que queremos estudiar en este trabajo.


\section{Clasificación de imágenes multi-etiqueta}

En cuanto al problema de clasificación multi-etiqueta en el reciente estudio Rethinking Multi-Label Image Classification With Deep Learning: Taxonomy, Challenge, and Outlook ~\cite{zhu2026rethinkingmultilabelimageclassification} analiza el cambio de paradigma en el problema que han traido los sistemas basados en redes convolucionales y transformers y estudia como diferentes métodos derivados de estas redes han mejorado la calidad de los sistemas de aprendizaje para esta tarea. Podemos encontrar un benchmarking de diferentes arquitecturas de redes como ResNet101 ~\cite{he2015deepresiduallearningimage}, ViT-Base ~\cite{dosovitskiy2021imageworth16x16words}, ML-GCN ~\cite{chen2019multilabelimagerecognitiongraph}, ASL ~\cite{benbaruch2021asymmetriclossmultilabelclassification}, Q2L ~\cite{liu2021query2labelsimpletransformerway}, SALGL ~\cite{10378641}, SGRE ~\cite{10546293}, and PAT ~\cite{xie2024counterfactualreasoningmultilabelimage}.

En cuanto a la resolución del problema conjunto de crear un clasificador multi-etiqueta dentro de un entorno federado, como ya hemos visto anterioremente hay estudios que se centran en evitar el problema haciendo un tratamiento previo de los datos para evitar el desbalanceo, otros estudian diferentes arquitecturas de la red como Transformer-based Federated Learning for Multi-Label Remote Sensing Image Classification ~\cite{2024transformerbasedfederatedlearningmultilabel}. Sin embargo estos estudios comparten la misma red entre los clientes o se aplican sobre datos tabulados. Nuestra propuesta requiere que la distribución de los datos contenga diferente número de etiquetas y que el sistema acepte una arquitectura con diferentes redes, no hemos encontrado ningún estudio que una todos estos conceptos en una solución.


%%
\cleardoublepage
